\chapter{The Message-Passing Framework}
\label{ch:message-passing}
\minitoc

A graph model must handle inputs of varying size, remain independent of 
node ordering, and incorporate graph structure into the computation. 
These requirements are captured by the relational inductive bias, together with the 
conditions of permutation invariance and equivariance. This chapter constructs 
a layer that satisfies these properties by design, assembling it from three 
component functions, examining the effects of stacking such layers, and 
interpreting the resulting formulation as a general design space for GNNs.

\section{From Relational Bias to Neighborhood Aggregation}
\label{sec:neighborhood-aggregation}

The relational inductive bias of Section~\ref{sec:relational-bias} asserts that
a node is characterized by its neighborhood. This fixes where a node's representation 
must come from: not from the node alone, but from the node together with 
the nodes adjacent to it. A layer that respects the bias therefore computes 
a node's representation by \textit{aggregating} the information 
held by its neighbors. This is the operative idea. The
question is how to perform that aggregation while meeting
the requirements the previous chapters imposed.

The CNN offers an instructive point of comparison,
because it faces the same problem on a fixed-size domain and solves it in a way
that does not transfer directly~\cite{bishop2024}. A convolutional layer reads
each position together with a fixed window of surrounding positions and applies
the same filter at every position. Two aspects of this design are worth
isolating. The window is of fixed size and fixed shape, so the filter always
reads the same number of inputs arranged in the same way; and the filter is
reused everywhere, so a single set of parameters serves the whole input. A
graph layer should follow the same instinct --- read a node together with its
local region rather than the whole graph --- but the neighborhood structure of
a graph differs fundamentally. A neighborhood is not a fixed window: different
nodes have different numbers of neighbors, so the operation must accept a
variable number of inputs. And a neighborhood has no fixed arrangement: it is a
set, with no first or last element, so the operation must not depend on the
order in which the neighbors are listed.

These two obstacles are met by two ingredients acting together, and the
distinction between what each contributes is worth drawing precisely, since it
is easy to attribute to one what belongs to the other. The first ingredient is
the aggregation itself. An operation that combines a neighborhood of arbitrary
size into a single vector resolves the problem of variable size directly: its
output has a fixed dimension regardless of how many neighbors entered it. If,
in addition, that operation is insensitive to the order of its inputs, it also
discards the arbitrary ordering, and the requirement of
Section~\ref{sec:permutation} is satisfied at the level of the neighborhood.
The aggregation, then, handles the variable size and supplies invariance to
the ordering of neighbors.

The second ingredient is \textit{parameter sharing}: the same functions, with
the same parameters, are applied at every node of the graph, exactly as a
single filter is reused at every position of an image. Its contribution is
distinct from that of the aggregation. Because the identical computation is
performed at every node, a relabeling of the nodes does not change the computation
applied to them. Combined with an order-insensitive aggregation, it allows the
desired permutation properties to be built into the layer.

Parameter sharing also keeps the number of parameters independent of the number
of nodes. The same layer can therefore be applied to graphs of arbitrary size,
which both improves scalability and permits a model trained on one graph to be
applied to another~\cite{bishop2024}. The two ingredients are therefore
complementary --- the parameter sharing guarantees that the same computation
is applied to every node, and the aggregation collapses a variable, unordered
neighborhood into a fixed vector --- and together they turn the relational inductive
bias into an operation that meets the requirements of the two preceding
chapters.

When a node's representation is formed by combining information from its neighborhood, 
through repeated aggregation and shared transformations, connected nodes tend to 
acquire increasingly similar representations. This behavior aligns the layer with
the \textit{homophily} assumption introduced in Section~\ref{sec:structural-properties}, namely 
that adjacent nodes are likely to be similar. For many graphs this bias is appropriate and is 
precisely what makes the approach effective. When neighboring nodes carry information 
that differs systematically from that of the central node, however, treating their contributions 
uniformly may become a limitation, a question taken up in Chapter~\ref{ch:limitations}.

\section{The Message-Passing Layer}
\label{sec:message-passing-layer}

The ideas developed in the previous section can now be stated formally. A layer of
the kind described takes the representations of all nodes and produces new
representations for each node, each formed from the node and its neighborhood.
This section assembles such a layer from three functions --- one that forms a
message from a neighbor, one that aggregates the messages of a neighborhood,
and one that updates a node from its aggregated neighborhood --- and states the
general equation they compose into. The unification of these three functions
into a single framework, under which many previously separate graph models
appear as instances, is due to Gilmer et al.~\cite{gilmer2017}; the exposition
here follows the formulation used throughout this work.

The object the layer transforms is the \textit{hidden representation} of a
node. Writing $\mathbf{h}_i^{(k)} \in \mathbb{R}^{d'}$ for the representation of
node $i$ after layer $k$, the layer maps the node representations
$\{\mathbf{h}_i^{(k-1)}\}$ to updated representations $\{\mathbf{h}_i^{(k)}\}$.
The representation at layer zero is initialized to the input features,
$\mathbf{h}_i^{(0)} = \mathbf{x}_i$, so the first layer reads the raw node
attributes and each subsequent layer reads the output of the one before.
The representations produced by the intermediate layers are therefore hidden. 
As with the feature matrix $\mathbf{X}$, whose row $i$ is the feature vector
$\mathbf{x}_i$, the hidden representations of all nodes are collected into a
matrix $\mathbf{H}^{(k)} \in \mathbb{R}^{n \times d'}$ whose row $i$ is
$\mathbf{h}_i^{(k)}$; the matrix form is used when a layer is written as a
single operation on the whole graph.

\subsection{The Message Function}
\label{sec:message-function}

The first function specifies what a neighbor contributes. For a node $i$ and a
neighbor $j \in \mathcal{N}(i)$, the \textit{message function} $\phi^{(k)}$
produces a vector --- the message sent from $j$ to $i$ --- from the
representations available at the two endpoints of the edge and, where present,
the feature of the edge itself:
\begin{equation}
  \mathbf{m}_{j \to i}^{(k)}
    = \phi^{(k)}\!\left(
        \mathbf{h}_i^{(k-1)},\,
        \mathbf{h}_j^{(k-1)},\,
        \mathbf{e}_{j,i}
      \right).
  \label{eq:message-function}
\end{equation}
The function $\phi^{(k)}$ is a learnable map whose output $\mathbf{m}_{j \to
i}^{(k)}$ is the message; its arguments are the representation of the receiving
node $i$, the representation of the sending neighbor $j$, and the feature
$\mathbf{e}_{j,i}$ of the edge from $j$ to $i$. Not every
argument need be used: a message may be formed from the neighbor alone, from
the neighbor together with the central node, or from either of these together
with the edge feature. The general form admits all three sources --- the
neighbor, the central node, and the edge --- and a particular architecture
selects among them. The edge feature is retained in the general form because 
the framework accommodates it, and because it is the point at which edge 
information would enter when available.

\subsection{The Aggregation Function}
\label{sec:aggregation-function}

The message function produces one message for each neighbor of $i$. Since a
node may have any number of neighbors, these messages form a collection of
variable size, and they must be combined into a single vector before the node
can be updated. The \textit{aggregation function} $\bigoplus$ performs this
combination over the neighborhood:
\begin{equation}
  \mathbf{a}_i^{(k)}
    = \bigoplus_{j \in \mathcal{N}(i)} \mathbf{m}_{j \to i}^{(k)}.
  \label{eq:aggregation-function}
\end{equation}
The aggregated vector $\mathbf{a}_i^{(k)}$ summarizes the entire neighborhood of
$i$ in a single vector of fixed dimension, whatever the number of neighbors,
which is what allows the layer to accept neighborhoods of varying size.
The neighborhood $\mathcal{N}(i)$ is a set, and
the messages carry no intrinsic order; the result of combining them must
therefore not depend on the order in which they are taken. The aggregation
function must be invariant to the ordering of its inputs. The operations
commonly used --- the sum, the mean, and the maximum, taken coordinatewise ---
all satisfy this condition, and any of them yields a well-defined aggregation.

\subsection{The Update Function}
\label{sec:update-function}

The aggregated neighborhood does not by itself determine the node's new
representation. Although information about the node itself may already enter
through the message function, the layer still requires a mechanism that
combines the neighborhood summary with the node's own prior representation. 
The \textit{update function} $\gamma^{(k)}$ combines the two, producing the new
representation from the node's previous representation and the aggregate of its
neighborhood:
\begin{equation}
  \mathbf{h}_i^{(k)}
    = \gamma^{(k)}\!\left(
        \mathbf{h}_i^{(k-1)},\,
        \mathbf{a}_i^{(k)}
      \right).
  \label{eq:update-function}
\end{equation}
The function $\gamma^{(k)}$ is a learnable map that takes the node's own prior
representation $\mathbf{h}_i^{(k-1)}$ together with the aggregated message
$\mathbf{a}_i^{(k)}$ and returns the updated representation $\mathbf{h}_i^{(k)}$.
Retaining the node's prior state alongside the neighborhood is what prevents the
node's own information from being discarded; the
balance between the two is a matter of how $\gamma^{(k)}$ is chosen, and
ranges from a simple sum of the two arguments to a learnable combination, a
question taken up with the design space in
Section~\ref{sec:design-space}.

\subsection{The General Equation of a Message Passing Layer}
\label{sec:general-equation}

Composing the three functions gives the general form of a message-passing
layer. Substituting the message function~\eqref{eq:message-function} into the
aggregation~\eqref{eq:aggregation-function}, and the result into the
update~\eqref{eq:update-function}, the representation of node $i$ at layer $k$
is
\begin{equation}
  \mathbf{h}_i^{(k)}
    = \gamma^{(k)}\!\left(
        \mathbf{h}_i^{(k-1)},\;
        \bigoplus_{j \in \mathcal{N}(i)}
          \phi^{(k)}\!\left(
            \mathbf{h}_i^{(k-1)},\,
            \mathbf{h}_j^{(k-1)},\,
            \mathbf{e}_{j,i}
          \right)
      \right).
  \label{eq:message-passing}
\end{equation}
This equation defines the message-passing framework (Figure~\ref{fig:message-passing-layer}), 
and the architectures of the following chapters are obtained by choosing the three
functions~$\phi^{(k)}$,~$\bigoplus$, and~$\gamma^{(k)}$. 
Where any of these functions applies 
a learnable linear transformation, the transformation is written with a weight matrix
$\mathbf{W}$, and an associated bias is always understood to be present even
where it is not written: the bias is frequently absorbed into the weight
notation or omitted for brevity, and its omission from a particular equation
never signifies its absence. Hamilton expresses the same construction with
operators named \textsc{aggregate} and \textsc{update}, in place of $\bigoplus$
and $\gamma^{(k)}$; the two formulations describe the same
computation~\cite{hamilton2020}.

The layer assembled in this way meets the permutation conditions of
Section~\ref{sec:permutation}: the message
function~\eqref{eq:message-function} is applied identically to every edge, the
aggregation~\eqref{eq:aggregation-function} is invariant to the order of the
neighbors, and the update~\eqref{eq:update-function} is applied identically to
every node. Relabeling the nodes by a permutation therefore does not change any
node's computed representation; it only moves that representation to the node's
new index. In the matrix notation of Section~\ref{sec:permutation}, a layer
that sends $(\mathbf{A}, \mathbf{H}^{(k-1)})$ to $\mathbf{H}^{(k)}$
sends the relabeled graph $(\mathbf{P}\mathbf{A}\mathbf{P}^\top,\,
\mathbf{P}\mathbf{H}^{(k-1)})$ to $\mathbf{P}\mathbf{H}^{(k)}$,
which is the condition of permutation equivariance
in~\eqref{eq:permutation-equivariance}. The per-node output follows the nodes
when they are reordered, as a node-level task requires. A task defined on the
whole graph requires instead a single output unchanged by relabeling --- the
permutation invariance of~\eqref{eq:permutation-invariance} --- and this is
obtained by reducing the per-node representations to one vector with a further
order-insensitive operation; that operation, called \textit{readout}, is the one 
used by graph-level architectures, and is developed in Chapter~\ref{ch:architectures}.

\begin{figure*}[t]
  \centering
  \includegraphics[width=\textwidth]{message-passing-layer}
  \caption[The three functions of a message-passing layer.]{%
    One message-passing layer applied to a central node $i$. For each neighbor
    $j \in \mathcal{N}(i)$, the message function $\phi^{(k)}$ forms a message
    from the representations of $j$ and $i$ and, where present, the edge
    feature. The aggregation $\bigoplus$ combines the neighbors' messages into a
    single vector, invariant to their order. The update function $\gamma^{(k)}$
    combines this aggregate with the node's own previous representation to
    produce $\mathbf{h}_i^{(k)}$.
  }
  \label{fig:message-passing-layer}
\end{figure*}

\section{Stacking Layers: Receptive Field and Depth}
\label{sec:receptive-field}

A single message-passing layer updates each node from its immediate neighbors.
Applying the layer once therefore lets each node incorporate information from 
one step away. To reach further, layers are stacked: the output representations of 
one layer become the input representations of the next, and a network of $K$ layers
applies the update~\eqref{eq:message-passing} $K$ times in succession, with
$k = 1, \ldots, K$. Each layer carries its own functions
$\phi^{(k)}, \gamma^{(k)}$ and its own parameters; the parameters are shared
across the nodes within a layer, as Section~\ref{sec:neighborhood-aggregation}
established, but not across layers, so the transformation applied at depth $k$
differs from the one applied at depth $k-1$.

Stacking has a definite effect on how far information can propagate into a 
node's representation, and the effect compounds with depth. After one layer,
$\mathbf{h}_i^{(1)}$ depends on the nodes in $\mathcal{N}(i)$, the neighbors at
distance one. After two layers, $\mathbf{h}_i^{(2)}$ is computed from the
first-layer representations of those neighbors, and each of those already
depends on its own neighbors; the representation of $i$ thus comes to depend on
the nodes within distance two. In general, after $k$ layers the representation
$\mathbf{h}_i^{(k)}$ depends on every node within geodesic distance $k$ of $i$
--- the set of nodes $j$ with $d(i, j) \le k$, the distance being the one
defined in Chapter~\ref{ch:graph-theory}. This is the \textit{receptive
field} of node $i$ at depth $k$: the region of the graph that can influence the
node's representation. Each additional layer expands the receptive field by one
\textit{hop}, so the depth of the network sets the radius of the neighborhood that each
node ultimately incorporates (Figure~\ref{fig:receptive-field}).

\begin{figure*}[t]
  \centering
  \includegraphics[width=\textwidth]{receptive-field}
  \caption[The receptive field of a node grows with depth.]{%
    Left: a target node in the input graph. Right: the computation that
    produces its representation, unrolled over message-passing layers. The
    first layer aggregates messages from the node's immediate neighbors; the
    second aggregates messages from the neighbors of those neighbors, and so
    on. After $k$ layers the representation of the target node depends on every
    node within $k$ hops of it, the receptive field at depth~$k$.
  }
  \label{fig:receptive-field}
\end{figure*}

The depth of the network thus controls the extent of the graph that each node's
representation can incorporate, and this is the quantity a designer sets in
choosing the number of layers. The final representations $\mathbf{h}_i^{(K)}$,
produced after the last of the $K$ layers, are the \textit{embeddings} of the
nodes: the learned vectors that summarize each node together with its
$K$-hop context, and that serve as the input to whatever prediction the task
requires --- a classifier applied per node, a score computed on pairs of nodes,
or a reduction over all nodes to a single vector for the graph.

\section{Transductive and Inductive Learning}
\label{sec:transductive-inductive}
 
The way a message-passing network is trained and evaluated depends on what
portion of the graph data is available during training. In practice, the two
task levels introduced in Chapter~\ref{ch:applications} commonly give rise to
two distinct learning settings. The distinction concerns what the network has 
access to during training and what is held back for evaluation, and it follows directly 
from the fact that a node's representation is computed from its neighborhood.
 
Consider first a task defined at the node level, in which each node carries a
label and the goal is to predict the labels of some nodes from those of others.
As in any supervised problem, the nodes are partitioned into a training set,
whose labels are used to fit the model, and an evaluation set, whose labels are
withheld and used to measure generalization. The graph structure, however,
prevents these two sets from being separated cleanly. When the network computes
the representation of a training node, it aggregates the representations of that
node's neighbors, and some of those neighbors belong to the evaluation set.
Their features and their position in the graph therefore enter the computation
performed during training, even though their labels do not. This is 
the \textit{transductive} setting. In node-classification problems it
is often described as a form of \textit{semi-supervised learning}, since the
labels of only a subset of the nodes are available during training.
 
A task defined at the graph level admits a cleaner separation. Here the dataset
is a collection of whole graphs, each carrying a single label, and the
partition into training and evaluation sets is a partition of the graphs. A
graph held out for evaluation is absent from training in its entirety --- its
nodes, its edges, and its label are all unseen --- because the unit of the
partition is the whole graph rather than a node within a shared structure. This
is the \textit{inductive} setting, and it is the one that arises in the usual
sense of training on some examples and generalizing to others never seen before.
 
What enables the inductive setting is the parameter sharing of
Section~\ref{sec:neighborhood-aggregation}. Because the functions
$\phi^{(k)}$ and $\gamma^{(k)}$ are shared across nodes and carry no dependence
on the identity or the number of nodes in any particular graph, a network fitted
on one set of graphs is immediately applicable to another: the same functions
compute representations on a graph the network has never encountered. 
The same property that decouples the parameter count from graph size 
is thus what allows a message-passing network to generalize across graphs.
 
\section{Expressive Power and the 1-WL Bound}
\label{sec:expressive-power}
 
The framework constructed so far guarantees that a message-passing layer
respects the permutation conditions of a graph, but it says nothing about how much 
the resulting representations can distinguish. A natural way to measure the
capability of a graph model is to ask which graphs it can tell apart: given two
graphs that are not isomorphic, does the network assign them different
representations, or does it map structurally different graphs to the same
output? This is the question of \textit{expressive power}, and it is the
property that governs how well a model can serve graph-level tasks, where
distinguishing one graph from another is the whole of what is required. Graph
isomorphism, defined in Chapter~\ref{ch:graph-theory}, is the precise notion of
two graphs being the same up to a relabeling; expressive power asks how close a
network comes to distinguishing every pair of graphs that are genuinely
different.
 
There is a classical procedure for distinguishing non-isomorphic graphs that
was introduced before GNNs and turns out to mirror message passing closely.
The one-dimensional Weisfeiler--Leman test, abbreviated \mbox{1-WL}, assigns a
color to every node and refines these colors iteratively. Initially all nodes
share a common color. At each round, a node's new color is determined by its
current color together with the collection of colors of its neighbors; nodes
that have the same color and whose neighbors have matching collections of colors
keep a common color, while any difference in the surrounding colors splits them
apart. The refinement continues until the partition of nodes into colors no
longer changes. Two graphs for which the refinement process produces 
different colorings are certainly not isomorphic.
 
The correspondence with message passing is structural and immediate. A round of
\mbox{1-WL} updates a node's color from its own color and the colors of its
neighbors; a message-passing layer updates a node's representation from
its own representation and the aggregated representations of its neighbors. The
color plays the role of the hidden representation, and the refinement step plays
the role of one layer. This parallel was established by
Morris et al.~\cite{morris2019}, and it has a consequence that bounds the
framework from above: a message-passing network can distinguish two graphs only
if \mbox{1-WL} assigns them different colorings. No network built from the
general equation~\eqref{eq:message-passing} can separate a pair of graphs that
\mbox{1-WL} fails to separate. The expressive power of message passing is thus bounded 
above by \mbox{1-WL}.
 
This bound is a ceiling: a particular architecture may fall short
of it, depending on how its aggregation and update functions are chosen.
Reaching the ceiling --- matching the full discriminative power of \mbox{1-WL}
rather than settling below it --- requires the aggregation to preserve enough
information about the neighborhood that distinct neighborhoods always yield
distinct aggregated vectors. This requirement strongly 
favors aggregation functions that preserve as much
neighborhood information as possible, and motivates the architecture developed
for graph-level tasks in Chapter~\ref{ch:architectures}. The question of what the
\mbox{1-WL} ceiling excludes --- the structurally different graphs that no
message-passing network can tell apart --- is a limitation of the framework and
is taken up in Chapter~\ref{ch:limitations}.
 
\section{The Design Space}
\label{sec:design-space}
 
The general equation~\eqref{eq:message-passing} is not a single model but a
template. It fixes the form of a layer --- a message formed per neighbor, an
aggregation over the neighborhood, an update of the node --- and requires only
that the three functions respect the permutation conditions; it does not fix
what those functions are. Every choice of $\phi^{(k)}$, $\bigoplus$, and
$\gamma^{(k)}$ that meets the conditions produces a valid layer, and the space
of such choices is the design space of the framework. The architectures studied
in this work are points in this space, and reading the equation as a space of
possibilities rather than a fixed formula is what makes their differences
understandable.
 
\subsection{Flexibility of the Framework}
\label{sec:framework-flexibility}
 
The flexibility of the equation can be seen function by function. The message
function $\phi^{(k)}$ may be as simple as a linear transformation of the
neighbor's representation, or as involved as an MLP applied to
the neighbor, the central node, and the edge feature together; it may apply one
weight matrix to the central node and a different one to its neighbors, allowing
a node to treat its own information differently from that of its neighborhood.
The aggregation $\bigoplus$ may be any order-insensitive operation, with the sum,
mean, and maximum being the common choices, each preserving a different
aspect of the neighborhood. The update $\gamma^{(k)}$ may be a linear map
followed by a nonlinearity, an MLP, or a rule that
concatenates the node's prior representation with the aggregated neighborhood,
$\mathbf{h}_i^{(k-1)} \oplus \mathbf{a}_i^{(k)}$, before transforming the
result, so that the two are kept distinct rather than summed. Each such decision
is a coordinate in the design space, and a complete architecture is a setting of
all of them.
 
The flexibility extends beyond the choice of functions to the very object whose
representation is learned. The equation forms a node's representation from
messages along its edges, but nothing in its structure requires the
unit of representation to be a node. A layer that built its messages from edge
features and aggregated over adjacent edges would, by the same mechanism, learn
representations of the edges themselves; the message, aggregation, and update
functions retain their roles, and only the object they act on changes. 
Edge-centered architectures found in the literature are instances of the
same framework, reoriented from nodes to edges. This generality indicates that
message passing is not tied to node representations specifically, but is a
template for computing on graph elements through local aggregation.
 
\subsection{The MLP as a Degenerate Case}
\label{sec:mlp-degenerate}
 
One point in the design space recovers a familiar model and locates the
framework relative to it. Suppose the message function is discarded and the
update depends on the node alone, so that the layer computes
$\mathbf{h}_i^{(k)} = \gamma^{(k)}(\mathbf{h}_i^{(k-1)})$ with no contribution
from the neighborhood. The graph structure plays no part; each node is
transformed in isolation, and the layer reduces to an MLP applied independently
to every node. The MLP is therefore the degenerate point of the design space at
which the neighborhood is ignored --- precisely the model that
Chapter~\ref{ch:nn-limitations} found unable to use graph structure.
 
Read in the other direction, this shows that a message-passing network is a
generalization of the MLP, with the MLP appearing as a particular instance of the framework. 
The structure that Chapter~\ref{ch:nn-limitations} identified as missing 
is supplied by exactly the terms that the degenerate case removes, which places 
the message-passing framework and the MLP at two ends of a single construction.
 
\subsection{Established Architectures as Instances}
\label{sec:established-architectures}
 
The design space is not merely a catalog of possibilities; it is populated by
architectures that have been proposed, studied, and applied, each fixing the
three functions in a particular way to serve a particular purpose. The
recognition that these architectures, developed separately, are instances of one
framework is itself the contribution of the message-passing
formulation~\cite{gilmer2017}: models that appear distinct on the surface differ
only in their choice of message, aggregation, and update.
 
This work examines three such instances, one representative of 
each task level introduced in Chapter~\ref{ch:applications}: a node-level architecture built on 
the assumption that adjacent nodes are similar and whose aggregation smooths 
each node toward its neighborhood; an edge-level architecture that uses node embeddings 
to score the plausibility of a link between two nodes; and a graph-level architecture whose 
aggregation is chosen to reach the \mbox{1-WL} ceiling of Section~\ref{sec:expressive-power}, 
maximizing the network's power to distinguish whole graphs. Each is the subject of the
 following chapter, where the general framework of this chapter is given a concrete form.
